% !PRINTED 144
\weberpage{144}
%<<CONT

worin die Coëfficienten $A_0,A_1,\ldots,A_m$ nur von $a$ und von den in $S$ etwa noch vorkommenden Coëfficienten abhängen. Wir bemerken aber ausdrücklich, dass, wenn in $S(\alpha,\beta)$ keine gebrochenen Zahlencoëfficienten vorkommen, auch in den Coëfficienten $A_0,A_1,\ldots,A_m$ keine Brüche auftreten.

Wir setzen nun
\begin{equation}
 \Phi(x)=A_0x^m+A_1x^{m-1}+\cdots+A_{m-1}x+A_m
\tag{4}
\end{equation}
dividiren $\Phi(x)$ durch $f(x)$ (nach \S.\,3) und erhalten einen Quotienten $Q$ und einen Rest, der in Bezug auf $x$ höchstens vom ersten Grade ist, also:
\begin{equation}
 \Phi(x)=Qf(x)+A+Bx;
\tag{5}
\end{equation}
hierin sind nun $A$ und $B$ ganze Functionen von $a$ und $b$, und auch sie enthalten keinerlei gebrochene Zahlencoëfficienten, wenn in $S$ keine solche vorkommen. Wenn wir nun $x=\alpha$ setzen, so ergiebt sich aus (5) und (3), da $f(\alpha)$ verschwindet,
\begin{equation}
 S(\alpha,\beta)=A+B\alpha.
\tag{6}
\end{equation}
Da aber $S(\alpha,\beta)$ und ebenso $A,B$ symmetrisch sind, so folgt durch Vertauschung von $\alpha$ und $\beta$
\begin{equation}
 S(\alpha,\beta)=A+B\beta.
\tag{7}
\end{equation}
Hieraus schliesst man, da $\alpha$ und $\beta$ von einander unabhängige Variable sind, dass $B=0$ und folglich
\begin{equation}
 S(\alpha,\beta)=A
\tag{8}
\end{equation}
sein muss, womit der Fundamentalsatz für diesen Fall bewiesen ist.

\begin{center}
{\large \S.\ 44.}

\vspace{0.5em}
{\large\gesperrt{Allgemeiner Beweis des Hauptsatzes.}}
\end{center}

Wir setzen nun voraus, der Fundamentalsatz sei bewiesen für symmetrische Functionen von \(n-1\) Veränderlichen und leiten ihn durch ein Verfahren, was dem in \S.\,43 angewandten ganz analog ist, für \(n\) Variable her.

Es sei wieder
\begin{equation}
 S=S(\alpha_1,\alpha_2,\ldots,\alpha_n)
\tag{1}
\end{equation}
eine ganze symmetrische Function der \(n\) Veränderlichen \(\alpha_1,\alpha_2,\ldots,\alpha_n\). Wenn wir sie nach Potenzen von \(\alpha_1\) ordnen und demgemäss setzen
\begin{equation}
 S=S_0\alpha_1^\mu+S_1\alpha_1^{\mu-1}+\cdots+S_{\mu-1}\alpha_1+S_\mu,
\tag{2}
\end{equation}
%CONT>>
